\chapter{Credit supply, refinancing and macro transmission}\label{ch:credit}
\section{A group-level causal product needs several levels of analysis}
Retail treatment effects do not exhaust the questions of a universal bank.
The institution must understand how funding conditions change credit supply,
how borrowers substitute between lenders, and how mortgage refinancing
transmits policy to household cash flow. Those questions use bank-firm
networks, market-level variation and aggregate time series. They also expose
why the product cannot apply one customer-level identification strategy
to every business problem.

Three published designs illustrate complementary information. The liquidity
shock study of \citet{mian2006} compares differently affected banks lending
to the same firm. \citet{agarwal2020} uses exposure to a refinancing
program to distinguish policy effects from borrowers' selected refinancing
choices. \citet{benetton2021} estimates mortgage demand and lender behavior
to study market-wide regulation. Each makes additional claims that require
additional assumptions as the analysis moves from local variation toward
aggregate counterfactuals.

\section{Separating loan supply from borrower demand}
Let $\Delta L_{fb}$ be loan growth for firm $f$ at bank $b$, and
let $Z_b$ measure an externally induced funding shock. A simple model is
\begin{equation}
 \Delta L_{fb}=\alpha_f+\beta Z_b+\epsilon_{fb}.
 \label{eq:bankfirm}
\end{equation}
The firm effect absorbs its common change in demand across lenders. For
firms with several bank relationships, subtracting the within-firm mean
gives
\[
 \Delta L_{fb}-\overline{\Delta L}_{f\cdot}
 =\beta(Z_b-\bar Z_{f\cdot})+
   (\epsilon_{fb}-\bar\epsilon_{f\cdot}).
\]
Identification comes from different lender shocks facing the same borrower,
under the restriction that their residual relationship-specific demand
components do not co-vary with those shocks. A firm's general investment
opportunity can correlate with the shocks of its chosen banks and still
be removed; a bank-specific export-credit need correlated with the funding
shock may remain a confounder.

The distinction is explicit in the archived 2006 paper
\citep[Section III.B, equation 5]{mian2006}. It also notes the importance
of unanticipated shocks and the restriction to multiple-bank borrowers.
The later question of total firm borrowing cannot simply insert the same
firm effect after aggregating to one firm-level observation. Borrower
substitution and aggregate real effects need additional evidence.

\section{An executed matched-lender study}
The synthetic construction has 600 firms, 24 banks and three relationships
per firm. Bank shocks $Z_b$ and firm demand terms $U_f$ are initially
independent standard normals. Relationship selection favors banks with
similar shock exposure through probabilities proportional to
$\exp(0.65U_f Z_b)$. This produces correlation between a firm's demand
and its observed lenders' shocks. Generated loan growth is
\begin{equation}
 \Delta L_{fb}=2-3Z_b+4U_f+\eta_b+\varepsilon_{fb},
 \label{eq:credit-dgp}
\end{equation}
with bank noise standard deviation $0.45$ and relationship noise standard
deviation $1.5$. Units are percentage points of loan growth per unit
funding shock. The true supply coefficient is $-3$ by construction.

Pooled regression leaves selected borrower demand in the error. Firm
demeaning removes the common $4U_f$ term. The study estimates both on
the same 1,800 relationships. The within estimator's illustrative interval
clusters residual score contributions by bank and uses a $t_{23}$
critical value. There are 24 independent bank shocks, not 1,800 independent
policy shocks. The demonstration is one sample under the stipulated law;
it does not qualify few-cluster inference for arbitrary lending networks.

\begin{table}[htbp]\centering\small
\begin{tabular}{@{}lrr@{}}
\toprule
Method & Coefficient & 95\% interval\\
\midrule
Generating effect & -3.000 & --\\
Pooled & -1.080 & --\\
Borrower effects & -2.923 & [-3.027, -2.819]\\
\bottomrule
\end{tabular}

\caption{Synthetic matched-bank loan-growth comparison. Relationship
selection biases the pooled association, while borrower effects recover
the stipulated supply coefficient much more closely.}
\label{tab:credit-supply}
\end{table}

The estimate changes from approximately $-1.08$ pooled to $-2.92$ with
borrower effects. The difference is economically substantial. It arises
because the design compares lenders serving the same borrower, rather
than because fixed effects are a universal solution to endogeneity.
If relationship-specific demand were correlated with $Z_b$, the within
estimator could remain biased. If all a firm's lenders had identical
exposure, that firm would contribute no within-firm identifying variation.

\section{What the bank would need to observe}
A single bank sees its own lending, not necessarily the same firm's other
relationships. It may need an appropriately permitted credit registry,
syndicated-loan dataset or borrower-supplied exposure record to implement
this design. A model for outside borrowing does not substitute for those
measurements without new assumptions. Entity resolution must preserve
legal borrower, consolidated group, facility and lender roles; several
facilities at one bank do not create several independent lender shocks.

The source snapshot must retain pre-shock relationships to avoid conditioning
on post-shock survival or new borrowing. Committed lines, drawn balances,
new originations, prices and collateral changes are different outcomes.
A borrower drawing an existing commitment during a shock can increase
observed balances while the bank tightens new credit. The product therefore
requires an outcome type before fitting the panel.

For a group credit committee, the first report can show the identified
relative supply response and its supported borrower population. A second
scenario may translate loan changes into funding, margin and capital
requirements. That translation needs the contract account, maturities,
loss assumptions and borrower substitution. The platform should preserve
the first empirical result even if the second remains model-dependent.

\section{Refinancing: eligibility, take-up and ownership}
Refinancing illustrates another frequent causal mistake: comparing borrowers
who refinance with those who do not. Refinancing choices respond to
attention, creditworthiness, equity, expected mobility and competing offers.
Those factors can also predict later spending and default. An observed
post-refinancing improvement does not by itself estimate the effect of
an offer or policy that induces refinancing.

\citet{agarwal2020} uses HARP eligibility to define treatment exposure,
with conforming high-LTV loans and a comparison group of observationally
similar non-agency loans. The paper examines pre-program comparisons,
refinancing, borrower savings and broader spending outcomes, and studies
competitive frictions in pass-through. The appropriate product lesson is
the separation of program exposure, take-up and downstream outcomes. Its
design and estimates depend on that institutional setting and should not
be transplanted unchanged to a new bank's retention program.

The bank also needs an ownership account. A lower mortgage rate can improve
the household's cash flow while reducing the value of a mortgage asset
held by the group. Origination or servicing fees may accrue to a different
entity. Paying off an old claim and originating a replacement cannot be
valued by adding the new loan's gross interest while ignoring the old
asset, funding and principal movements. The causal module supplies a
supported behavioral contrast; the mortgage account supplies the signed
cash and continuation value for each owner.

A practical study can define assigned refinancing assistance as treatment,
measure application, offer, acceptance and closing as separate stages,
and value the original eligible cohort. Failed delivery and non-take-up
remain in an intention-to-treat analysis. A complier effect or mechanism
analysis needs additional assumptions. Outcome maturity can be long even
when take-up is observed quickly, which affects both study power and the
product investment calendar.

\section{Why equilibrium counterfactuals need more structure}
A broad change in capital requirements or a group-wide mortgage price
policy can alter competitor prices, market shares and loan amounts.
\citet{benetton2021} models borrowers' discrete product choice and
continuous loan demand, together with multiproduct lender pricing. It
uses variation in predetermined risk-weighted capital requirements as
cost shifters and separates lender and product variation. This construction
allows equilibrium questions beyond a local reduced-form effect, under
the demand, supply, exclusion and behavioral restrictions of that model.

For the proposed platform, this is a reason to expose a structural extension
interface. A fitted local response may be an input or a validation target
for a structural model, but it does not determine the entire market
counterfactual. Supply response, outside options and competitive conduct
remain separate parameters or assumptions. Good market-share prediction
under the incumbent is informative about fit and incomplete evidence about
an unprecedented policy intervention.

\citet{beraja2019} makes the aggregation issue particularly relevant to
monetary transmission: regional variation can reveal heterogeneous responses
while common aggregate effects cancel from the comparison. Adding local
coefficients does not recover what was differenced out. A group-level
report should distinguish directly identified local effects, reweighted
population effects and equilibrium model scenarios.

\section{Treasury transmission and the shock catalogue}
The time-series workspace should preserve a catalogue of candidate shocks:
their economic definition, observation timestamp, normalization, source,
instrument construction and maintained exclusion restrictions. A market
rate change, a policy-announcement surprise and a bank-specific liquidity
shock are different objects. A discovered lead-lag relation can suggest
a candidate mechanism, but it does not turn an ordinary residual into an
identified structural shock.

For a selected shock, local projections and structural VARs can be
challengers with a common outcome, horizon and normalization. The external-
instrument argument in \citet{stock2018} requires the relevant dynamic
exclusions; \citet{jorda2024} explains estimation and inference alternatives.
The report should show response covariance, weak-instrument diagnostics,
regime sensitivity and the difference between a shock response and a
sustained policy path. Few independent macro episodes remain a limitation
even in a bank with billions of transaction records.

\section{What the shared product contributes}
These applications share temporal reconstruction, economic entity roles,
source lineage, score covariance and evidence export. They retain distinct
identification routes: matched borrowers for supply, policy eligibility
for refinancing, structural demand and supply for equilibrium, and external
shock restrictions for macro transmission. This is a stronger reusable
product than a collection of regressions with common plotting defaults.

MacroFinance can consume a supported shock or empirical response while
retaining its structural equilibrium assumptions. SMEwallet can consume
an observed or identified relationship response while preserving uncertainty
about outside borrowing. FinMarketing can use an eligible-cohort response
in a compatible mortgage or relationship account. The proposal specifies
those exchanges; it does not certify current adapters or endorse a copied
historical repository milestone as a current qualification.
